% Source: https://artofproblemsolving.com/community/c3h2573303_algebraic_manipulation
% Stable identity: aops-topic-2573303-q01
Put $t=\sqrt{b-\sqrt b}$. For $b=1$ we get $t=0$ and $x=1$, so the expression is undefined; hence $b\ge2$ and $0<t<b$. Then
\[x+1=\dfrac{2b}{b-t},\qquad x-1=\dfrac{2t}{b-t},\]
so
\[a=\dfrac{(x+1)^4}{x(x-1)^2}=\dfrac{\dfrac{16b^4}{(b-t)^4}}{\dfrac{b+t}{b-t}\cdot\dfrac{4t^2}{(b-t)^2}}=\dfrac{4b^4}{t^2\left(b^2-t^2\right)}.\]
With $s=\sqrt b$ we have $t^2=b-s$ and $b^2-t^2=b^2-b+s$. Using $s^2=b$, the denominator is
\[(b-s)\left(b^2-b+s\right)=b^3-b^2-b+\left(2b-b^2\right)s.\]

\textbf{Case 1: $b$ is not a perfect square.} Then $s$ is irrational. Since $a$ is rational, so is the denominator $\dfrac{4b^4}a$, which forces $2b-b^2=0$, i.e. $b=2$. Then the denominator is $8-4-2=2$ and $a=\dfrac{4\cdot16}{2}=32$.

\textbf{Case 2: $b=s^2$ with an integer $s\ge2$.} The denominator is $\left(s^2-s\right)\left(s^4-s^2+s\right)=s^2(s-1)\left(s^3-s+1\right)$, so
\[a=\dfrac{4s^6}{(s-1)\left(s^3-s+1\right)}.\]
Both $s-1$ and $s^3-s+1$ are coprime to $s$, so each must divide $4$. From $s-1\mid4$ we get $s\in\{2,3,5\}$, but then $s^3-s+1=7,25,121$, none of which divides $4$. No solutions.

So the only pair is $\boxed{(a,b)=(32,2)}$.
